
\subsection{\texorpdfstring{\textbf{
      Реализация скоростного градиента}}
  {Реализация скоростного градиента}}

Спустя множество рассуждений было решено использовать мультиагентный подход,
когда множество агентов - целевых состояний изменяются по одной и той же динамике
непрерывной сети Хопфилда, а настройка происходит с использованием того факта,
что они должны быть стационарными точками.

Так появляется следующая математическая модель:

\[
  \frac{dx_{i}^\mu}{dt}= \sum\limits_{i=1}^N W_{ij}y_{j}^\mu
  - \beta x_{i}^\mu + b_{i}
\]

\[
  y_{j}^\mu = \tanh{(x_{j}^\mu)},
  \quad i,j=1,\dots, N, \quad \mu=1,\dots, P
\]

Здесь \(P\) - число целевых паттернов, а \(N\) - число нейронов


Берётся интегральный функционал
\[
  Q(t) = \int R(y(s),W(s),b(s),s) \, ds
\]

\[
  R(y(t),W(t),b(t),t) = \frac{1}{2P}\sum\limits_{\mu=1}^P
  {\lVert r^\mu(t) \rVert^2}
\]

\[
  r^\mu(t)=-x_{*}^\mu+W(t)y_{*}^\mu(t)+b(t)
\]

\[
  x_{*_{i}}^\mu = atanh{((1-\epsilon)\xi_{i}^\mu)}
\]

где \( \xi_{i}^\mu \) - целевые паттерны для обучения.

Тогда

\[
  \frac{dW}{dt}= -\gamma_{W}\P_{S}(\nabla_{W}{R(\tilde{y},W,b,t)}),
  \quad W(0)\in S
\]

\[
  \frac{db}{dt}= -\gamma_{b}\nabla_{b}R(\tilde{y},W,b,t)
\]

\[
  \tilde{y}_{i}^\mu(t) = y_{i}^\mu(t) + d_{i}^\mu \\
  \quad d_{i}^\mu = A_{i}^\mu(e^{-\alpha t} + delta)\sin{\omega_{i}^\mu t}
\]

На изображениях можно видеть явное улучшение по сравнению с Хеббом.
Стационарные паттерны остаются после 1000 шагов (Рис.
\ref{media:image_velgrad_stationary}).
\begin{figure}[h]
  \centering
  \includegraphics[width=0.5\textwidth]{image_velgrad_stationary.png}
  \caption{}
  \label{media:image_velgrad_stationary}
\end{figure}

Уже после первого шага убирается практически весь шум (Рис.
\ref{media:image_velgrad}).

\begin{figure}[h]
  \centering
  \includegraphics[width=0.4\textwidth]{image_velgrad.png}
  \caption{}
  \label{media:image_velgrad}
\end{figure}

\begin{figure}[h]
  \centering
  \includegraphics[width=0.5\textwidth]{image_velgrad_2.png}
  \caption{}
  \label{media:image_velgrad_2}
\end{figure}

После 5 шагов изображения полностью восстановлены (Рис.
\ref{media:image_velgrad_2}).